% H15: semantic information through natural scrambles. Sources: hypotheses/H15-semantic-information-scrambles/{README.md (Round 1b), summary/content.tex, summary/meta.json};
% writeup/visuals/H44-context-wipe/{caption.tex, README.md} (shared H44+H15 figure); interpretation/swarm-constants.json (ΔW_erase).
\section{H15: Memory loss costs nothing visible; a context wipe costs a tenth of output}

\textbf{The question:}
Which stored information keeps an agent working: its memory file, its context window, the chat or its artifacts?
We cannot run new swarms, so we use natural scrambles from the logs.
These are memory losses, newcomers with empty memories, chat cuts, artifact switches and forced context wipes.

\textbf{The intuition:}
Kolchinsky--Wolpert semantic information is the information whose loss lowers a system's viability, a measure of how well it keeps going.
The test is an intervention on an information channel: scramble a store and watch whether output falls.
An agent has two stores with different timescales.
Memory holds information over days; the context window holds it over one session of about 40 calls.
A scramble erases one store, or cuts the channel from the room to the agent.

\textbf{What we measured:}
At the day scale, the observable is a viability $V^*$ per agent-day, chosen by a pre-registered procedure (D2.6).
Regime I uses engagement and regime III uses reliability (one minus the real-failure share).
Each event is compared with placebo days of agents in the same unit, with a Kalman-filtered counterfactual.
At the call scale, the observable is work commits per call around a forced wipe at the 41-turn cap.
It compares calls $+1\ldots+10$ with calls $-20\ldots-11$, on the DQ1 ledger and the DQ4 work ledger.
The scope is goal periods outside the reserved data.
\emph{Scheduler field:} removed by same-day differencing against other agents, and by the comparison within one segment.
\emph{External drive:} partly removed; the cap times forced wipes, but 118/172 artifact switches fall on goal-change days.
The pre-registered forced-versus-voluntary contrast also carries a task-phase confound.
\emph{Shared model prior} and \emph{common input:} not relevant; the designs are within one agent and make no influence claim.
\emph{Synthetic check:} 150 replicates on the real village skeleton.
The calibrated null gives 1--7\% false positives, but the day-scale power is only 17--41\%.
The check chose the Kalman counterfactual because it removes the bias from resets that happen on bad days.

\begin{figure}[h]
\includegraphics[width=\textwidth]{H44-context-wipe/fig.pdf}
\caption{The forced context wipe, shared with H44.
Writes and work commits dip right after the wipe and recover over a few calls, against a no-reset null that keeps rising.
Panel (e) carries H15's ledger estimate: work commits fall 39\% for ten calls in 7/9 periods.
Animation: \texttt{writeup/visuals/H44-context-wipe/anim.mp4}.}
\label{fig:int-H15}
\end{figure}

\textbf{What we found:}
Memory losses show no day-scale cost.
Losing 50--80\% of memory (12 events) gives $+0.29$\,SD [$-0.05$, $+0.63$] on $V^*$.
Newcomers with empty memories give $+0.20$ [$-0.13$, $+0.54$], and they progress \emph{more} ($+0.31$ [0.06, 0.56]).
Mid-period artifact switches cost nothing measurable on work output ($-0.00$ [$-0.17$, $+0.17$]).
The context window is the one store whose loss shows.
A forced wipe cuts work commits by 39\% [$-42$, $-35$] over ten calls, in 7/9 periods.
Write calls fall 26\% (8/9) and real failures rise 27\%.
The lost work is 2.5--13\% of a forced segment's output, with a median near 10\%.
The memory written at the wipe does not soften the dip ($\rho=+0.01$; top against bottom tercile $-0.05$ [$-0.10$, $+0.01$]).
H44 measures the same dip ($-38\%$) with its own pipeline.
Scope: regime III for the wipe; all regimes for the day-scale nulls.
The claim stands\cred{0.78}: ``a forced wipe cuts work commits by about 39\% for ten calls. More memory does not help, and memory carries no day-scale information.''
Its expected usefulness is EU~2.73 (value if true 3.5; two raters, one blind).
Small conflict: \texttt{meta.json} counts ``turns''; the card's round 1b counts calls on the ledger, and this note follows the card.

\textbf{What it means:}
For a scaffold builder: each forced wipe at the 41-turn cap costs about 39\% of committed work for ten calls, about a tenth of output.
To recover it, raise the cap or carry a task-state summary across the reset; nobody has tested the summary yet.
Do not ask the agent to write more memory at the wipe; the dose has no effect.
Do not spend effort to protect memory files against loss: no day-scale cost appears, though the tests can only see large effects.

\textbf{Caveats:}
The day-scale tests have low power (17--41\%), so ``no cost'' means ``no large cost''.
The viability measure $V^*$ changed between rounds, so the period verdicts are replications under a new yardstick.
The wipe dip is post hoc: the pre-registered forced-versus-voluntary test came out with the wrong sign.
NE41 then replicated the dip with dated predictions on new timing and output data.
The study measures neither the semantic information $S$ nor its share $\eta$; Kolchinsky--Wolpert supplies the vocabulary only.
The chat-cut test has 3 events.
The fragile flag is set: the coordinator rated the claim fragile.
The confirmation script must be re-frozen on work commits.
No confirmation test on reserved goal periods has run.
